\chapter{一般の分布}
    \section{\texorpdfstring{$\bm{x}$}{x}の分散共分散行列を\texorpdfstring{$V=\Var{}{\bm{x}}$}{V=Var(x)}とすると\texorpdfstring{$\Var{}{\bm{v}^\top\bm{x}} = \bm{v}^\top V\bm{v}$}{Var(v\string^T x) = v\string^T V v}}
        \begin{proof}
        \quad\par
            $\hat{\bm{x}} = \E{}{\bm{x}}$とすると
            \begin{align*}
                \Var{}{\bm{v}^\top\bm{x}} &= \E{}{\left(\bm{v}^\top\bm{x}-\E{}{\bm{v}^\top\bm{x}}\right)^2} = \E{}{\left(\bm{v}^\top\bm{x}-\bm{v}^\top\hat{\bm{x}}\right)^2} = \E{}{\left(\bm{v}^\top(\bm{x}-\hat{\bm{x}})\right)^2}\\
                &= \E{}{\bm{v}^\top\left(\bm{x}-\hat{\bm{x}}\right)\left(\bm{x}-\hat{\bm{x}}\right)^\top\bm{v}} = \bm{v}^\top\E{}{\left(\bm{x}-\hat{\bm{x}}\right)\left(\bm{x}-\hat{\bm{x}}\right)^\top}\bm{v} = \bm{v}^\top V\bm{v}
            \end{align*}
        \end{proof}
    \section{ベクトル確率変数の線形写像の像の期待値と分散共分散行列}
        \begin{shadebox}
            $X_1,\dotsc,X_n\in\realNumbers$を確率変数、$A\in\realNumbers^{n\times m}$、$\bm{X}\coloneqq[X_1,\dotsc,X_n]^\top$、$\Sigma_X\coloneqq\CovMat{\bm{X}}$とすると次が成り立つ。
            \[\E{}{A^\top\bm{X}} = A^\top\E{}{\bm{X}},\quad\CovMat{A^\top\bm{X}} = A^\top\Sigma_X A\]
            特に$X_1,\dotsc,X_n$が独立で同一の分散$\sigma^2$をもち、かつ$A$が$n$次直交行列であれば
            \[\CovMat{A^\top\bm{X}} = \Sigma_X\]
        \end{shadebox}
        \begin{proof}
            \newcommand{\ai}{{\bm{a}_i}}
            \newcommand{\aj}{{\bm{a}_j}}
            \quad\par
            $\E{}{A^\top\bm{X}} = A^\top\E{}{\bm{X}}$は期待値の線形性より明らか。$\CovMat{A^\top\bm{X}} = A^\top\Sigma_X A$を示す。
            $\bm{Y}\coloneqq A^\top\bm{X}$とし、$A$の第$i$列ベクトルを$\bm{a}_i$とすると
            \begin{align*}
                \CovMat{A^\top\bm{X}}_{ij} &= \CovMat{\bm{Y}}_{ij} = \Cov{Y_i}{Y_j} = \E{}{\left(Y_i-\E{}{Y_i}\right)\left(Y_j-\E{}{Y_j}\right) } \\
                &= \E{}{\left(\ai^\top\bm{X}-\E{}{\ai^\top\bm{X}}\right) \left(\aj^\top\bm{X}-\E{}{\aj^\top\bm{X}}\right)} \\
                &= \E{}{\ai^\top\left(\bm{X}-\E{}{\bm{X}}\right)\aj^\top\left(\bm{X}-\E{}{\bm{X}}\right)} \\
                &= \E{}{\ai^\top\left(\bm{X}-\E{}{\bm{X}}\right)\left(\bm{X}-\E{}{\bm{X}}\right)^\top\aj} \\
                &= \ai^\top\E{}{\left(\bm{X}-\E{}{\bm{X}}\right)\left(\bm{X}-\E{}{\bm{X}}\right)}\aj = \ai^\top\Sigma_X\aj
            \end{align*}
            これより$\CovMat{A^\top\bm{X}} = A^\top\Sigma_X A$。
            \par
            特に$X_1,\dotsc,X_n$が独立で同一の分散$\sigma^2$をもち、かつ$A$が$n$次直交行列であれば
            \[\CovMat{A^\top\bm{X}} = A^\top\Sigma_X A = A^\top\sigma^2I_n A = \sigma^2I_n = \Sigma_X\]
        \end{proof}
    \section{スコア関数}
        $\theta$をパラメータとする確率密度関数$f(x;\theta)$に従うi.i.d確率変数$X_1,\dotsc,X_n \in \Omega$に対して「スコア関数」を次式で定義する。
        \[ V_n(\bm{X},\theta) \coloneqq \partDerivLong{\log f(\bm{X};\theta)}{\theta}{} = \partDerivLong{\log \prod_{i=1}^n f(X_i;\theta)}{\theta}{} = \partDerivLong{\sum_{i=1}^n \log f(X_i;\theta)}{\theta}{} = \sum_{i=1}^n V(X_i,\theta) \]
        \subsection{\texorpdfstring{$\E{\bm{X}}{V(\bm{X},\theta)} = 0$}{E\_X (V(X,Θ)) = 0}}
            \begin{proof}
                \begin{align*}
                    \E{\bm{X}}{V_n(\bm{X},\theta)} &= \integrate{\Omega^n}{}{V_n(\bm{x},\theta) \prod_{i=1}^n f(x_i;\theta)}{}{\bm{x}} = \integrate{\Omega^n}{}{\partDerivLong{\prod_{i=1}^n f(x_i;\theta)}{\theta}{}}{}{\bm{x}} \\
                    &= \partDerivLong{\integrate{\Omega^n}{}{\prod_{i=1}^n f(x_i;\theta)}{}{\bm{x}}}{\theta}{} \quad (\text{微分と積分の順序交換可能性を仮定}) \\
                    &= \partDeriv{1}{\theta}{} = 0
                \end{align*}
            \end{proof}
    \section{Cram\'er-Raoの不等式}
        \begin{shadebox}
            $\theta$をパラメータとする確率密度関数$f(x;\theta)$に従うi.i.d確率変数$X_1,\dotsc,X_n \in \Omega$から計算した$\theta$の不偏推定量$\hat{\theta}(\bm{X})$について次式が成り立つ
            \[ \Var{\bm{X}}{\hat{\theta}(\bm{X})} \geq \left(n\E{\textcolor{red}{X}}{V(\textcolor{red}{X},\theta)^2}\right)^{-1} = \left(n\Var{X}{V(X,\theta)}\right)^{-1} \]
        \end{shadebox}
        \begin{proof}
            \quad\par
            (参考: \url{https://to-kei.net/estimator/cramer-rao/})\\
            \begin{align*}
                \Var{\bm{X}}{\hat{\theta}(\bm{X})}\E{\bm{X}}{V_n(\bm{X},\theta)^2} &= \integrate{\Omega}{}{\left(\hat{\theta}(\bm{x}) - \theta\right)^2 f(\bm{x};\theta)}{}{\bm{x}} \integrate{\Omega}{}{V_n(\bm{x},\theta)^2 f(\bm{x};\theta)}{}{\bm{x}} \\
                &= \integrate{\Omega}{}{\left(\left(\hat{\theta}(\bm{x}) - \theta\right)\sqrt{f(\bm{x};\theta)}\right)^2}{}{\bm{x}} \integrate{\Omega}{}{\left(V_n(\bm{x},\theta)\sqrt{f(\bm{x};\theta)}\right)^2}{}{\bm{x}} \\
                &\geq \integrate{\Omega}{}{\left(\hat{\theta}(\bm{x}) - \theta\right)V_n(\bm{x},\theta)f(\bm{x};\theta)}{}{\bm{x}} \quad (\text{Cauchy-Schwarzの不等式を用いた}) \\
                &= \integrate{\Omega}{}{\hat{\theta}(\bm{x})V_n(\bm{x},\theta)f(\bm{x},\theta)}{}{\bm{x}} \quad \left(\because\; \text{直前の定理より }\E{\bm{X}}{V_n(\bm{X},\theta)} = 0\right) \\
                &= \integrate{\Omega}{}{\hat{\theta}(\bm{x})\left(\partDerivLong{\log f(\bm{x};\theta)}{\theta}{}\right)f(\bm{x};\theta)}{}{\bm{x}} = \integrate{\Omega}{}{\hat{\theta}(\bm{x})\partDerivLong{f(\bm{x};\theta)}{\theta}{}}{}{\bm{x}} \\
                &= \partDerivLong{\integrate{\Omega}{}{\hat{\theta}(\bm{x})f(\bm{x};\theta)}{}{\bm{x}}}{\theta}{} \quad (\text{微分と積分の順序交換可能性を仮定した}) \\
                &= \partDerivLong{\E{\bm{X}}{\hat{\theta}(\bm{X})}}{\theta}{} = \partDeriv{\theta}{\theta}{} = 1
            \end{align*}
            ここで
            \begin{align*}
                \E{\bm{X}}{V_n(\bm{X},\theta)^2} &= \E{\bm{X}}{\left(\sum_{i=1}^n V(X_i,\theta)\right)^2} = \sum_{i=1}^n\sum_{j=1}^n\E{\bm{X}}{V(X_i,\theta)V(X_j,\theta)} \\
                &= \sum_{i=1}^n \E{\bm{X}}{V(X_i,\theta)^2} \quad (\because\; V(X_i,\theta),V(X_j,\theta),\;i\neq j\text{は独立で}\E{\bm{X}}{V(X_i,\theta)} = 0) \\
                &= \sum_{i=1}^n \E{\textcolor{red}{X_i}}{V(X_i,\theta)^2} = n\E{X}{V(X,\theta)^2}
            \end{align*}
            以上より定理の主張が成り立つ。
        \end{proof}
    \section{条件付き確率}
        \subsection{\texorpdfstring{$\Pr{A}=\sum_k{\Pr{A,B_k}} = \sum_k{\cPr{A}{B_k}\Pr{B_k}}$}{Pr(A) = Σ\_k Pr(A, Bk) = Σ\_k Pr(A|Bk)Pr(Bk)}}
        \label{tech:確率の分解}
        \subsection{\texorpdfstring{$\cPr{A,B}{C}=\cPr{A}{B,C}\cPr{B}{C}$}{Pr(A,B|C) = Pr(A|B,C)Pr(B|C)} (スライディング)}
            \begin{proof}
                \[\cPr{A,B}{C} = \frac{\Pr{A,B,C}}{\Pr{C}} = \frac{\cPr{A}{B,C}\Pr{B,C}}{\Pr{C}} = \cPr{A}{B,C}\cPr{B}{C}\]
            \end{proof}
        \subsection{\texorpdfstring{$\cPr{A}{C} = \sum_k{\cPr{A}{B_k,C}\cPr{B_k}{C}}$}{Pr(A|C) = Σ\_k Pr(A|Bk,C)Pr(Bk|C)}}
            \begin{proof}
                \begin{align*}
                    \cPr{A}{C} &= \frac{\Pr{A,C}}{\Pr{C}} = \frac{\sum_k{\Pr{A,B_k,C}}}{\Pr{C}} \quad(\because\;\ref{tech:確率の分解}) \\
                    &= \frac{\sum_k{\cPr{A}{B_k,C}\Pr{B_k,C}}}{\Pr{C}} = \sum_k{\cPr{A}{B_k,C}\cPr{B_k}{C}}
                \end{align*}
            \end{proof}
    \section{条件付き確率密度関数}
        連続型2変数$X,Y$に関して、条件付き確率密度関数を次のように定義する。
        \[p_Y(y|X=x) \coloneqq \frac{p(x,y)}{\integrate{\ymin}{\ymax}{p(x,y)}{}{y}}\]
        ($p(x,y)$は$X,Y$の同時確率密度関数)
    \section{条件付き期待値}
        2変数$X,Y$に関して、条件付き期待値を次のように定義する。
        \begin{align*}
            \textrm{(離散)} &: \E{Y}{Y|X=x} \coloneqq \sum_y y\Pr{Y=y|X=x} = \sum_y y\frac{\Pr{Y=y,X=x}}{\Pr{X=x}} \\
            \textrm{(連続)} &: \E{Y}{Y|X=x} \coloneqq \integrate{\ymin}{\ymax}{yp_Y(y|X=x)}{}{y}
        \end{align*}
        \subsection{\texorpdfstring{$\E{X,Y}{\E{Y}{Y|X}} = \E{X,Y}{Y}$}{E\_(X,Y) (E\_Y (Y|X)) = E\_(X,Y) (Y)}}
            \begin{shadebox}
                条件付き期待値について次の性質が成り立つ。
                \[\E{X,Y}{\E{Y}{Y|X}} = \E{X,Y}{Y}\]
            \end{shadebox}
            \begin{proof}
                \quad\\
                (離散型)
                \begin{align*}
                    \E{X,Y}{\E{Y}{Y|X}} &= \sum_x\sum_y\E{Y}{Y|X=x}\Pr{X=x,Y=y} \\
                    &= \sum_x\E{Y}{Y|X=x}\sum_y\Pr{X=x,Y=y} \quad (\because\;\E{Y}{Y|X=x}\textrm{は}y\textrm{に無関係}) \\
                    &= \sum_x\left(\sum_y{\frac{\Pr{X=x,Y=y}}{\Pr{X=x}}}\right)\Pr{X=x} \\
                    &= \sum_x\sum_y y\Pr{X=x,Y=y} = \E{X,Y}{Y}
                \end{align*}
                (連続型)\\
                $X,Y$の同時確率密度関数を$p(x,y)$とし、$X$の周辺確率密度関数を$p_X(x)$とする。
                \begin{align*}
                    \E{X,Y}{\E{Y}{Y|X}} &= \integrate{\xmin}{\xmax}{\integrate{\ymin}{\ymax}{\E{Y}{Y|X=x}p(x,y)}{}{y}}{}{x} \\
                    &= \integrate{\xmin}{\xmax}{\E{Y}{Y|X=x}\integrate{\ymin}{\ymax}{p(x,y)}{}{y}}{}{x} \quad (\because\;\E{Y}{Y|X=x}\textrm{は}y\textrm{に無関係}) \\
                    &= \integrate{\xmin}{\xmax}{\left(\integrate{\ymin}{\ymax}{y\frac{p(x,y)}{p_X(x)}}{}{y}\right)p_X(x)}{}{x} = \integrate{\xmin}{\xmax}{\integrate{\ymin}{\ymax}{yp(x,y)}{}{y}}{}{x} = \E{X,Y}{Y}
                \end{align*}
            \end{proof}
    % \begin{comment}
    % \section{指示関数}
    %     \subsection{同時事象の指示関数の期待値}
    %         \begin{shadebox}
    %             同時事象の指示関数の期待値に関して次が成り立つ。
    %             \begin{align*}
    %                 \E{}{\ind{Y \satisfy C_Y,\;X \satisfy C_X}} &= \E{X}{\ind{X\satisfy C_X}\E{Y}{\cind{Y\satisfy C_Y}{X\satisfy C_X}}} \\
    %                 &= \E{X}{\ind{X\satisfy C_X}\cPr{Y\satisfy C_Y}{X\satisfy C_X}}
    %             \end{align*}
    %         \end{shadebox}
    %         \begin{proof}
    %             \begin{align*}
    %                 \textrm{左辺} &= \E{}{\cind{Y\satisfy C_Y}{X\satisfy C_X}\ind{X\satisfy C_X}} \\
    %                 &= \sum_x{\sum_y{\Pr{X=x,Y=y}\cind{y\satisfy C_Y}{x\satisfy C_X}\ind{x\satisfy C_X}}} \\
    %                 &= \sum_x{\sum_y{\cPr{Y=y}{X=x}\Pr{X=x}\cind{y\satisfy C_Y}{x\satisfy C_X}\ind{x\satisfy C_X}}} \\
    %                 &= \sum_x{\Pr{X=x}\ind{x\satisfy C_X}\sum_y{\cPr{Y=y}{X=x}\cind{y\satisfy C_Y}{x\satisfy C_X}}} \\
    %                 &= \sum_x{\Pr{X=x}\ind{x\satisfy C_X}\E{Y}{\cind{Y\satisfy C_Y}{x\satisfy C_X}}} \\
    %                 &= \E{X}{\ind{X\satisfy C_X}\E{Y}{\cind{Y\satisfy C_Y}{X\satisfy C_X}}} \\
    %                 &= \E{X}{\ind{X\satisfy C_X}\cPr{Y\satisfy C_Y}{X\satisfy C_X}}
    %             \end{align*}
    %         \end{proof}
    % \end{comment}
    \section{確率不等式}
        \subsection{Markovの不等式}
            \begin{shadebox}
                $(X,\mathcal{M},\mu)$を測度空間、$f:X\to[-\infty,\infty]$を可測関数、$a>0$とするとき、次が成り立つ。
                \[\displaystyle \mu(\{x\in X\,|\,|f(x)|\ge a\})\le\frac{1}{a}\int_{X}|f(x)|\mu(dx)\]
            \end{shadebox}
            \begin{proof}
                \quad\par
                (引用 : \url{http://mathcommunication.hatenablog.com/entry/2017/06/27/001722})
                \par
                \begin{equation*}
                    \begin{split}
                        \int_{X}|f(x)|\mu(dx) & =\int_{|f(x)|\ge a} |f(x)|\mu(dx) + \int_{|f(x)|<a}|f(x)|\mu(dx) \\
                        & \ge\int_{|f(x)|\ge a} |f(x)|\mu(dx) \\
                        & \ge a\,\mu(\{x\in X\,|\,|f(x)|\ge a\})
                    \end{split}
                \end{equation*}
            \end{proof}
        \subsection{Chebyshevの不等式}
            \begin{shadebox}
                任意の連続型確変数$X$について次が成り立つ。
                \[ \Pr{|X-\mu|\geq a\sigma} \leq 1/a^2 \]
                但し$a>0, \mu = \E{}{X}, \sigma^2 = \Var{}{}{X}$
            \end{shadebox}
            \begin{proof}
                \quad\par
                $X$の確率密度関数を$f(x)$とすると
                \begin{align*}
                    \Pr{|X-\mu| \geq a\sigma} &= \left(\int_{-\infty}^{\mu-a\sigma} + \int_{\mu+a\sigma}^\infty\right) \leq \left(\int_{-\infty}^{\mu-a\sigma} + \int_{\mu+a\sigma}^\infty\right)\frac{(x-\mu)^2}{a^2\sigma^2}f(x)\mathrm{d}x \\
                    &\left(\because\;\frac{(x-\mu)^2}{a^2\sigma^2} \geq 1, x\in(-\infty,\mu-a\sigma]\cup[\mu+a\sigma,\infty)\right) \\
                    &\leq \integrate{-\infty}{\infty}{\frac{(x-\mu)^2}{a^2\sigma^2}f(x)}{}{x} = 1/a^2
                \end{align*}
            \end{proof}
        \subsection{H\"{o}ffdingの補題}
            \begin{shadebox}
                $(\Omega,\mathcal{F},P)$を確率空間とし、$X$を$E_{P}[X]=0$、$a\leq X\leq b$ (a.e.)を満たす確率変数とする。このとき任意の$t>0$に対して次の不等式が成り立つ。
                \[E_{P}[\exp(tX)]\le\exp\left(\frac{1}{8}t^{2}(b-a)^{2}\right)\]
            \end{shadebox}
            \begin{proof}
                \quad\par
                (参考 : \url{http://mathcommunication.hatenablog.com/entry/2017/06/27/001722})
                \par
                $E_{P}[X]=0$より$a<0,b>0$であることに注意。指数関数の凸性より、
                \[\exp(tX)\le\frac{X-a}{b-a}\exp(tb)+\frac{b-X}{b-a}\exp(ta),\quad(\mathrm{a.e.})\]
                であるから
                \begin{equation}
                    E_{P}[\exp(tX)]\le\frac{-a}{b-a}\exp(tb)+\frac{b}{b-a}\exp(ta)
                \end{equation}
                今、
                \[\displaystyle p\coloneqq\frac{-a}{b-a},\quad u\coloneqq t(b-a)\]
                とおくと$\displaystyle p\coloneqq\frac{-a}{b-a},\quad u\coloneqq t(b-a)$であり、
                \begin{align*}
                    (1) &= p\exp(tb)+(1-p)\exp(ta)\\
                    &= p\exp((1-p)u)+(1-p)\exp(-pu)
                \end{align*}
                \begin{screen}
                    $\because\;\;1-p=\frac{b}{b-a}$より$tb = t(b-a)(1-p) = (1-p)u$ また、$ta = -tp(b-a) = -pu$
                \end{screen}
                さらに、$\phi(u) \coloneqq \log{\left(\textrm{式}(1)\right)}$とおくと
                \begin{align*}
                    \phi(u) &= \log\left[p\exp( (1-p)u)+(1-p)\exp(-pu)\right]\\
                    &= \log\left[\exp( (1-p)u)\cdot(p+(1-p)\exp(-u))\right]\\
                    &= (1-p)u+\log(p+(1-p)\exp(-u))
                \end{align*}
                \begin{gather*}
                    \phi'(u)=1-p+\frac{-(1-p)\exp(-u)}{p+(1-p)\exp(-u)},\\
                    \phi''(u)=\frac{p(1-p)\exp(-u)}{(p+(1-p)\exp(-u))^{2}}
                \end{gather*}
                ここでTaylorの定理より、ある$u_0\in(0,u)$が存在して
                \[\phi(u)=\phi(0)+\phi'(0)u+\frac{1}{2}\phi''(u_{0})u^{2}=\frac{1}{2}\phi''(u_{0})u^{2}\]
                であるが、相加・相乗平均の関係より
                \[\phi''(u_{0})=\frac{p\cdot(1-p)\exp(-u_{0})}{(p+(1-p)\exp(-u_{0}))^{2}}\le\frac{1}{4}\]
                であるので
                \[\phi(u)\le\frac{1}{8}u^{2}\]
                よって
                \[E_{P}[\exp(tX)]\le\exp(\phi(u) )\le\exp\left(\frac{1}{8}u^{2}\right)=\exp\left(\frac{1}{8}t^{2}(b-a)^{2}\right)\]
            \end{proof}
        \subsection{H\"{o}ffdingの不等式}
            \begin{shadebox}
                $(\Omega,\mathcal{F},P)$を確率空間とし、$X_{1},\ldots,X_{n}$を独立同分布な確率変数で$a_{i}\le X_{i} \le b_{i}\ (\mathrm{a.e.})$を満たすとする。このとき任意の$\epsilon>0$に対して次の不等式が成り立つ。
                \begin{gather*}
                    P(S-E_{P}[S]\ge\varepsilon)\le \exp\left(-\frac{2\varepsilon^{2}}{\sum_{i=1}^{n}(b_{i}-a_{i})^{2}}\right),\\
                    P(S-E_{P}[S]\le-\varepsilon)\le\exp\left(-\frac{2\varepsilon^{2}}{\sum_{i=1}^{n}(b_{i}-a_{i})^{2}}\right)
                \end{gather*}
            \end{shadebox}
            \begin{proof}
                \quad\par
                (引用 : \url{http://mathcommunication.hatenablog.com/entry/2017/06/27/001722})
                \par
                任意の$t>0$に対して
                \begin{align*}
                    P(S-E_{P}[S]\ge\varepsilon) &= P(\exp(t(S-E_{P}[S]))\ge\exp(t\varepsilon))\\
                    &\le \exp(-t\varepsilon)E_{P}[\exp(t(S-E_{P}[S]))] & \text{（マルコフの不等式）}\\
                    &= \exp(-t\varepsilon)\prod_{i=1}^{n}E_{P}[\exp(t(X_{i}-E_{P}[X_{i}]))] & (\text{各}X_i\text{は独立})\\
                    &\leq \exp(-t\varepsilon)\prod_{i=1}^{n}\exp\left(\frac{t^2}{8}(b_i-a_i)^2\right) & \text{(ヘフディングの補題)}\\
                    &= \exp(-t\varepsilon)\exp\left(\frac{t^2}{8}\sum_{i=1}^{n}(b_{i}-a_{i})^{2}\right)\\
                    &= \exp\left(\frac{t^2}{8}\sum_{i=1}^{n}(b_{i}-a_{i})^{2} - t\varepsilon\right)
                \end{align*}
                が成り立つ。
                \[\frac{1}{8}t^{2}\sum_{i=1}^{n} (b_{i}-a_{i})^{2}-t\varepsilon = \frac{1}{8}\left(\sum_{i=1}^{n}(b_{i}-a_{i})^{2}\right)\left(t-\frac{4\varepsilon}{\sum_{i=1}^{n}(b_{i}-a_{i})^{2}}\right)^{2}-\frac{2\varepsilon^{2}}{\sum_{i=1}^{n}(b_{i}-a_{i})^{2}}\]
                であり、第1項$\ge0$だから
                \[P(S-E_{P}[S]\ge\varepsilon)\le\exp\left(-\frac{2\varepsilon^{2}}{\sum_{i=1}^{n}(b_{i}-a_{i})^{2}}\right)\]
                同様にして
                \[\displaystyle P(S-E_{P}[S]\le-\varepsilon)\le\exp\left(-\frac{2\varepsilon^{2}}{\sum_{i=1}^{n}(b_{i}-a_{i})^{2}}\right)\]
            \end{proof}
    \section{中心極限定理}
        \renewcommand{\A}{\frac{\sqrt{n}(\hat{\mu}_n-\mu)}{\sigma}}
        \newcommand{\SUM}{\sum_{j=1}^n}
        \newcommand{\PI}{\prod_{j=1}^n}
        \begin{shadebox}
            標準化された標本平均$\A$の分布は標準正規分布に弱収束する。すなわち
            \[\lim_{n\rightarrow\infty}\Pr{\A\leq x} = \Phi(x)\]
            但し$\Phi(x)$は標準正規分布の累積分布関数である。
        \end{shadebox}
        \begin{proof}
            \quad\par
            (引用 : \url{https://ja.wikipedia.org/wiki/中心極限定理})
            \par
            \[Z_n \coloneqq \A = \frac{\SUM{X_j}-n\mu}{\sqrt{n}\sigma} = \SUM{Y_j/sqrt{n}}\quad\textrm{where }Y_j \coloneqq \frac{X_j-\mu}{\sigma}\]
            とおく。$Y_j$の平均は0分散は1である。
            \par
            $Z_n$の特性関数$\phi_{Z_n}(t)$は
            \begin{align*}
                \phi_{Z_n}(t) &= \E{}{e^{itZ_n}} = \E{}{e^{it\SUM{Y_j/\sqrt{n}}}} = \E{}{\PI{e^{itY_j/\sqrt{n}}}} = \PI{\E{}{e^{itY_j/\sqrt{n}}}}\quad(\because\textrm{各}Y_j\textrm{は独立})\\
                &= \PI{\phi_{Y_j}(t/\sqrt{n})} = \left[\phi_{Y_1}(t/\sqrt{n})\right]^n\quad(\because\textrm{各}Y_j\textrm{は全部同等})
            \end{align*}
            $\phi_{Y_1}(t)$のTaylor展開を考える。
            \begin{align*}
                \phi_{Y_1}(0) &= 1\\
                {\phi_{Y_1}}'(0) &= \left.\E{}{iY_1e^{itY_i}}\right|_{t=0} = 0\\
                {\phi_{Y_1}}''(0) &= \left.\E{}{-{Y_1}^2e^{itY_i}}\right|_{t=0} = -\E{}{{Y_1}^2} = -1
            \end{align*}
            これより、
            \[\phi_{Y_1}(t/\sqrt{n}) = 1 - \frac{1}{2}\frac{t^2}{n} + c\frac{1}{3!}\left(\frac{t}{\sqrt{n}}\right)^3 + o\left(\left(\frac{t}{\sqrt{n}}\right)^3\right)\quad(t\rightarrow0)\]
            よって
            \[\phi_{Z_n}(t) = \left[1 - \frac{1}{n}\frac{t^2}{2} + c\frac{1}{3!}\left(\frac{t}{\sqrt{n}}\right)^3 + o\left(\left(\frac{t}{\sqrt{n}}\right)^3\right)\right]^n \rightarrow e^{-t^2/2}\quad\textrm{as }n\rightarrow\infty\]
            最右辺は標準正規分布の特性関数であり、特性関数と確率密度関数は1対1で対応するから$Z_n = \A$は$n\rightarrow\infty$で標準正規分布に従う。
        \end{proof}